\chapter{Analisi}

\section{Requisiti}

Il progetto mira alla realizzazione di un videogioco educativo chiamato
``//TODO'', che consente al giocatore di familiarizzarsi con i linguaggi
di programmazione della famiglia Assembler. La programmazione Assembler
è caratterizzata dall'esecuzione sequenziale di istruzioni semplici e
dal flusso di controllo gestito tramite salti condizionali.

\subsection{Requisiti funzionali}

\begin{itemize}
\item
  Il videogioco deve presentare al giocatore una serie di puzzle da
  risolvere, ognuno con istruzioni chiare su cosa implementare e
  strumenti per verificare che il codice scritto sia corretto ed
  efficiente.
\item
  Il videogioco deve fornire all'utente un ambiente di sviluppo semplice
  e intuitivo, basato sul drag\&drop. L'ambiente deve permettere sia
  l'aggiunta, modifica e rimozione delle istruzioni, sia la possibilità
  di annullare le modifiche recenti per tornare al codice scritto in
  precedenza.
\item
  Durante l'esecuzione, il videogioco deve rappresentare graficamente
  ogni operazione eseguita, in modo che il giocatore possa comprendere
  cosa stia facendo il proprio programma.
\item
  ll videogioco deve permettere sia l'esportazione che l'importazione
  del programma in un formato testuale facilmente modificabile con un
  editor di testo, per condividere il proprio programma con altri
  giocatori o per provare soluzioni altrui.
\item
  Il videogioco deve fornire gli strumenti necessari per risolvere bug
  nel codice dell'utente, come per esempio la possibilità di eseguire il
  programma step by step.
\end{itemize}

\subsection{Requisiti non funzionali}

\begin{itemize}
\item
  Far sì che la VM non sia bloccante (i.e.~la sua esecuzione non deve
  rallentare o bloccare l'interazione con l'utente o il gioco intero).
\end{itemize}

\section{Analisi e modello del dominio}

Il videogioco deve contenere al suo interno una virtual machine in grado
di eseguire le istruzioni fornite dall'utente e dare indicazioni circa
eventuali errori presenti nel codice. La lista di istruzioni (il
``programma'') deve tenere traccia di tutte le più recenti modifiche
effettuate, per poterle annullare o eventualmente riapplicare al
programma. Ogni livello dovrà fornire all'utente un sottoinsieme di
tutte le possibili istruzioni eseguibili dalla VM per la risoluzione del
particolare problema che propone; dovrà inoltre essere possibile
verificare la correttezza dell'algoritmo proposto dall'utente e
accertarsi che tale soluzione non funzioni in maniera fortuita per un
solo particolare input. Ciò richiederà di avere una entità in grado di
confrontare l'output generato dal programma dell'utente con quello
fornito dalla soluzione ``canonica'' su più possibili sequenze di input
e non sulla singola sequenza proposta all'utente. Tale approccio renderà
poi necessario dover mostrare all'utente eventuali input per i quali il
programma sviluppato non funzioni.

\begin{figure}[H]
\includegraphics[width=\textwidth]{images/analysis/domain}
\caption{Schema UML rappresentante le principali entità prese in
    considerazione nella descrizione del dominio. \texttt{ExecutionController}
rappresenta l'entità preposta a controllare la correttezza del programma
dell'utente (il quale viene associato a un particolare livello tramite
l'entità \texttt{LevelController})}
\end{figure}
